\section{CONCLUSION AND FUTURE WORK}\label{sec:conclusion}

We built and evaluated a policy-gated multi-agent design for cloud pipeline governance against a
live remote language model, explicit non-agent baselines, an adversarial evaluation against an
independent oracle, and closed-form sensitivity analyses. The design only partly holds up. Against
static orchestration it improves recovery time and cost by smaller margins than the original report
claimed, it increases manual interventions, and two idealized non-LLM comparators recover faster
still. The most consequential result, though, concerns the evaluation itself: run through the same
code path, the identical configuration reaches very different conclusions under a deterministic mock
than under a live model, by roughly three orders of magnitude on recovery time and by a comparably
large margin on decision quality. For this class of system, \emph{which} evaluation path produced a
number matters at least as much as the number itself. The practical takeaway is narrow: this is not a
verdict on LLM agents in operations generally, but a reason to ask what evaluated a given governance
claim, and to treat a mock number and a live number as separate findings until both exist.

Future work splits into two parts. One is the evidence this campaign did not collect: a live
cross-model sweep to test model-agnosticism, an exercise of the trust core's mode transitions and
kill-switch latency under load, and replacing the two modelled baselines with deployed controllers
and real production traces (Section~\ref{sec:discussion}). The other is design work: replacing the
static agent-priority bid with an explicit negotiation protocol, switching on the
bounded-adaptation mechanism currently left off to preserve determinism, and extending the evaluation
to multi-cloud deployments and formal verification of the policy layer.
